\lysh{20941}{2}
\begin{alist}
\item Έχουμε:
\[P(-2) = (-2)^3 + 2(-2)^2 + (-2) + 3 = 1.\]
Επομένως το $-2$ δεν είναι ρίζα του πολυωνύμου.
\item Το σχήμα Horner με διαιρετέο $P(x)$ και διαιρέτη $x + 2$ δίνει
\begin{center}
\Horner[quotient=false]{1,2,1,3}{-2}
\end{center}
Επομένως το πηλίκο της διαίρεσης είναι το πολυώνυμο
\[\pi(x) = x^2 + 1.\]
\item Η ταυτότητα της διαίρεσης $P(x) : (x + 2)$ γράφεται
\[P(x) = (x^2 + 1)(x + 2) + 1.\]
\end{alist}

